% ============================================================
% ЗАКЛЮЧЕНИЕ
% ============================================================
\section*{ЗАКЛЮЧЕНИЕ}
\addcontentsline{toc}{section}{ЗАКЛЮЧЕНИЕ}

В ходе выполнения выпускной квалификационной работы была спроектирована и реализована интерактивная обучающая веб-платформа SopromatLab, предназначенная для изучения основ сопротивления материалов с мгновенной визуализацией напряжённо-деформированного состояния балочных конструкций.

В процессе работы выполнены следующие задачи:

\begin{enumerate}
  \item \textbf{Формализованы математические модели} для трёх типов статически определимых конструкций: консольной балки, шарнирно-опёртой балки и балки с вылетом. Для каждого типа определены граничные условия, уравнения равновесия и формулы определения опорных реакций.

  \item \textbf{Разработаны алгоритмы численного расчёта} эпюр изгибающих моментов~$M(x)$, поперечных сил~$Q(x)$, продольных сил~$N(x)$ и прогибов~$v(x)$. Расчёт основан на методе сечений (ROSE) с суперпозицией вкладов нагрузок и численном интегрировании методом трапеций. Погрешность определения прогибов не превышает~0{,}5\% при 500~точках дискретизации.


  \item \textbf{Спроектирована архитектура приложения} по методологии Feature-Sliced Design, обеспечивающая модульность и масштабируемость. Реализовано шесть архитектурных слоёв: app, pages, features, entities, widgets, shared.

  \item \textbf{Реализован интерактивный редактор балок} с визуализацией деформированной формы в реальном времени. Время от изменения параметра до обновления визуализации составляет менее~20~мс. Визуализация включает цветовую дифференциацию зон растяжения и сжатия, интерактивный hover с отображением значений в произвольной точке.

  \item \textbf{Созданы образовательные модули:} раздел теории (8~разделов), пошаговое решение (6~шагов), контекстные подсказки (21~условие), тренажёр с генерацией случайных задач и автоматической проверкой ответов, 5~учебных сценариев из области транспортной инженерии.

  \item \textbf{Реализованы функции экспорта:} генерация PDF-отчёта формата A4 с эпюрами, характерными значениями; обмен конфигурацией по ссылке через Base64-кодирование.

  \item \textbf{Интегрирован контекстный AI-консультант на базе большой языковой модели} (LLM): реализован виджет ChatAssistant, передающий текущее состояние расчётной модели в системный промпт по парадигме retrieval-augmented generation. Интеграция выполнена через единый шлюз OpenRouter, что позволяет переключаться между провайдерами без изменения клиентского кода. Консультант объясняет физический смысл эпюр и помогает пользователю интерпретировать результаты расчёта.
\end{enumerate}

Платформа реализована как полностью клиентское приложение (React~19 + TypeScript~5.9 + Vite), не требующее серверной инфраструктуры для расчётов. Все вычисления выполняются в браузере пользователя, что обеспечивает мгновенный отклик и возможность работы без подключения к интернету.

Разработанная платформа отличается от существующих онлайн-калькуляторов тем, что изначально проектировалась как интерактивный учебный продукт с педагогическим слоем: контекстными подсказками, пошаговым решением, тренажёром и учебными сценариями.

\textbf{Научная и практическая значимость.} Научная новизна работы заключается в комплексном объединении в одной образовательной системе интерактивного редактора расчётной схемы, контекстной системы подсказок, режима тренажёрной самопроверки и AI-консультанта на базе большой языковой модели, получающего полное состояние расчётной модели через RAG-инъекцию контекста. Практическая значимость состоит в создании готового к использованию инструмента, который может быть применён в учебном процессе по дисциплинам <<Сопротивление материалов>>, <<Механика деформируемого твёрдого тела>>, <<Прикладная механика>> в технических вузах. Контрольные расчёты для четырёх классических задач показали отклонение прогибов от аналитических решений менее~0{,}1\%, что подтверждает корректность реализованного расчётного ядра.

\textbf{Перспективы развития:}
\begin{itemize}
  \item расширение на статически неопределимые системы (метод сил, метод перемещений);
  \item добавление рамных конструкций (Г-образная, П-образная и многопролётные рамы);
  \item переход на дообученную (fine-tuned) языковую модель, специализированную на инженерной механике, и расширение базы знаний AI-консультанта учебниками по сопромату;
  \item локализация интерфейса на английский язык;
  \item разработка мобильной версии с адаптивным интерфейсом;
  \item реализация модуля сложного сопротивления (косой изгиб, кручение) и расчёта на устойчивость.
\end{itemize}

Таким образом, поставленные задачи выпускной квалификационной работы выполнены в полном объёме. Разработанная веб-платформа SopromatLab представляет собой работоспособный, точный и педагогически обоснованный инструмент для интерактивного обучения основам сопротивления материалов и может быть рекомендована к использованию в учебном процессе технических вузов.





